% Covariance-control result: independently audited PASS; see critical-constant-rate-fresh-audit.txt.
\section{Bounded critical rate for all functions at quadratic depth}
\label{sec:critical-constant-rate}

Covariance changes are cheaper than coherent mean shifts when the
relative change is small. The following sharper capacity obstruction
uses that fact. It concerns the critical variational rate, and does not
compare the optimal rates of competing functions closely enough to
prove a posterior sample-complexity lower bound.

\begin{theorem}[Universal bounded rate at quadratic depth]
\label{thm:critical-constant-rate}
There are universal finite constants $A,K_0>0$ such that, for every
$d\ge1$, there is an integer depth $L_d\le A(d+1)^2$ for which
\[
 J_{L_d}(f)\le K_0
 \qquad\hbox{for every Boolean }f:\{-1,1\}^d\to\{-1,1\}.
\]
The depth $L_d$ is the same for all functions on that cube.
\end{theorem}

We first record exact scalar formulas for a controlled feature law.
Put
\[
 F(q)=\mathbb E\tanh^2(\sqrt q Z),\qquad
 a(q)=1-F(q),\qquad
 R(q)=\frac{q[1-F(q)]^2}{F(q)}\quad(q>0).
\]
The value $R(0)=1$ makes it continuous at zero.

\begin{lemma}[Rank-one covariance control]
\label{lem:rankone-critical-control}
Suppose a feature covariance $Q$ has every diagonal entry equal to $q$
and $Q\succeq vff^T$ for $v>0$, where $f$ is Boolean.
For any $\epsilon\ge0$, let the next feature law be coordinatewise tanh
of a centered Gaussian with covariance $Q+\epsilon vff^T$.
Its entropy relative to the ordinary next-layer law $\mathsf M_Q$ is
at most $\epsilon^2/4$. Its second moment $Q'$ has constant diagonal
$q'$ and obeys $Q'\succeq v'ff^T$, where
\begin{equation}\label{eq:critical-covariance-control}
 u=q+\epsilon v,\qquad
 q'=F(u),\qquad v'=(1+\epsilon)v[1-F(u)]^2.
\end{equation}
Consequently, for $w=v/q$ and $w'=v'/q'$,
\[
 w'=w\frac{1+\epsilon}{1+\epsilon w}R(u).
\]
All these scalar quantities depend only on $q,v,\epsilon$, not on the
other entries of $Q$ or on the choice of $f$.
\end{lemma}

\begin{proof}
The Loewner bound implies $f\in\operatorname{range}(Q)$ and
$v f^TQ^+f\le1$. The two Gaussian laws have the same support, and the
rank-one determinant identity gives their relative entropy as
\[
 \frac12\{t-\log(1+t)\},\qquad
 t=\epsilon v f^TQ^+f\le\epsilon.
\]
This is at most $t^2/4\le\epsilon^2/4$. Applying coordinatewise tanh
can only decrease relative entropy.
Write the changed Gaussian as
$\sqrt{(1+\epsilon)v}\,fZ+G$, where $Z$ is standard Gaussian and
$G$ is an independent centered Gaussian of covariance $Q-vff^T$.
Its conditional feature mean given $Z$ is $f\psi(Z)$, since the
variances of the coordinates of $G$ all equal $q-v$.
Thus $Q'\succeq\mathbb E\psi(Z)^2 ff^T$.
Gaussian integration by parts and projection onto the first Hermite
polynomial give
\[
 \mathbb E\psi(Z)^2\ge(\mathbb E Z\psi(Z))^2
     =(1+\epsilon)v\,[\mathbb E\operatorname{sech}^2(\sqrt u Z)]^2.
\]
This is the claimed $v'$, and the diagonal formula and ratio identity
follow directly. In particular, these formulas automatically retain
$0<v'\le q'$, since they arise from a valid covariance lower bound.
\end{proof}

\begin{proof}[Proof of Theorem~\ref{thm:critical-constant-rate}]
Write $D=d+1$ and $m=2^d$. Begin with the unperturbed first feature law.
Its covariance $Q$ has constant diagonal
$q_0=F(D/d)\in[F(1),F(2)]$, and
Lemma~\ref{lem:first-tanh-spectrum} implies
\[
 Q\succeq e^{-20D}I\succeq v_0ff^T,
 \qquad v_0=e^{-20D}/m.
\]
This first layer costs zero entropy. We will use
Lemma~\ref{lem:rankone-critical-control} at each subsequent layer,
tracking its guaranteed scalar component $v$, not the largest possible
rank-one component of the covariance.

We collect universal scalar estimates. The function $F$ is increasing,
$0<F(q)<q$ for $q>0$, and Taylor expansion gives
\[
 F(q)=q-2q^2+O(q^3),\qquad R(q)=1+O(q^2).
\]
There consequently exist universal $u_0>0$ and $C\ge1$ such that,
for $0<u\le u_0$,
\begin{equation}\label{eq:critical-local-estimates}
 F(u)\le u-u^2,\qquad R(u)\ge1-Cu^2.
\end{equation}
We may and do take $u_0\le1/4$.
A useful global lower bound is
\begin{equation}\label{eq:critical-variance-lower}
 F(q)\ge\frac q{1+3q}.
\end{equation}
Indeed $\tanh^2 t\ge t^2/(1+t^2)$ follows from $|\sinh t|\ge|t|$.
Cauchy--Schwarz applied to $V=qZ^2$ gives
$\mathbb E[V/(1+V)]\ge(\mathbb EV)^2/\mathbb E[V(1+V)]
=q/(1+3q)$.

Choose a sufficiently small universal $v_*>0$ satisfying
\[
 v_*\le\min\{u_0/4,1/64,1/(100C)\}.
\]
Then choose a positive universal $c$ so small that
\[
 c\le\min\{v_*/4,u_0/3,1/3,1/(144C)\}.
\]
Put $\delta=c/D$.

\emph{Free cooling.}
Apply the unmodified next-layer law, that is, $\epsilon=0$, until
$q\le\delta$ for the first time. This costs no entropy.
There is a universal $c_F>0$ with
$q-F(q)\ge c_Fq^2$ for $0<q\le F(2)$: the ratio is positive and
continuous, and tends to two at zero. Hence the number of cooling
layers is at most a universal constant times $1/\delta$.
The final cooled variance $q_c$ lies in $[\delta/2,\delta]$, by
monotonicity and~\eqref{eq:critical-variance-lower}.

During cooling the ratio $w=v/q$ is multiplied by $R(q)$ at each step.
This loses at most a fixed positive factor, uniformly in the number
of cooling layers. To see this directly, on $[0,F(2)]$ one has
$|\log R(q)|\le C_R q^2$ for a universal $C_R$, by continuity,
positivity and the expansion at zero. Also
$\sum q_s^2\le c_F^{-1}\sum(q_s-q_{s+1})\le F(2)/c_F$.
Thus the cooled ratio satisfies
\[
 w_c\ge c_0v_0,
 \qquad \log(1/w_c)\le A_0D
\]
for universal $c_0>0,A_0<\infty$.

\emph{Amplifying the relative component.}
Now set $\epsilon=\delta$ at each layer until $w\ge1/2$.
Throughout this phase $q\le2\delta$. Indeed
$u=q+\delta v\le(1+\delta)q$, and the increasing function $u-u^2$
on the relevant interval gives
\[
 q'\le2\delta(1+\delta)-[2\delta(1+\delta)]^2\le2\delta.
\]
While $w\le1/2$, the exact ratio identity and
$u\le3\delta$ imply
\[
 \log(w'/w)
 \ge\frac\delta4-18C\delta^2\ge\frac\delta8.
\]
Here we used $\log(1-Cu^2)\ge-2Cu^2$, valid by the chosen constants.
The number $\tau$ of these layers is therefore at most
$8A_0D/\delta+1=O(D^2)$. Their total entropy is at most
\[
 \tau\delta^2/4\le2A_0D\delta+\delta^2/4=O(1).
\]
Furthermore, since $u\ge q$, monotonicity and
\eqref{eq:critical-variance-lower} imply
$1/q'\le1/q+3$ at every layer. Thus, at the end of this phase,
\[
 q\ge\frac1{2/\delta+3\tau}\ge\frac{c_1}{D^2},
 \qquad v\ge\frac q2\ge\frac{c_1}{2D^2}
\]
for some universal $c_1>0$.
At the same time $v\le q\le2\delta<v_*$.

\emph{Growing the absolute component.}
Set $\epsilon=16v$ until $v\ge v_*$.
The invariant $w\ge1/3$, or equivalently $q\le3v$, is preserved.
To check this, while $v<v_*$ one has
$u=q+16v^2\le4v$. If $w\ge1/2$, the ratio identity gives
$w'\ge\frac12(1-16Cv^2)\ge1/3$.
If $1/3\le w\le1/2$, it gives
\[
 \log(w'/w)\ge4v-32Cv^2\ge0.
\]
The initial ratio is at least $1/2$, so induction proves the invariant.
Using $F(u)\le u$ also gives
\[
 v'\ge v(1+16v)(1-2u)
       \ge v(1+16v)(1-8v)\ge v+4v^2,
 \qquad v'\le v+16v^2.
\]
Thus $1/v-1/v'\ge2$ until the stopping time, which is $O(D^2)$.
The final value is between $v_*$ and $2v_*$, and the total entropy
of this phase telescopes:
\[
 \sum\frac{(16v)^2}{4}
   =64\sum v^2\le16\sum(v'-v)\le32v_*.
\]

Finally choose output $z=f$. Its signed margin is one, and
$Q\succeq v_*ff^T$ bounds its Gaussian readout cost by $1/(2v_*)$.
All entropy costs plus this readout cost are bounded by a universal
constant $K_0$.
The number of hidden layers, and therefore the depth, is at most a
universal constant times $D^2$.
Crucially, $q_0,v_0$, every subsequent scalar recursion, every control,
and every stopping rule depend only on $d$. Although the full
covariance matrices depend on $f$, the resulting depth $L_d$ is common
to all Boolean functions. This proves the theorem.
\end{proof}

This construction sharpens an absolute-rate capacity obstruction. It
still does not show that the optimal rate of a complex competitor is
below, or sufficiently close to, the optimal rate of a specified simple
target. That comparison remains essential for a posterior lower bound.
